\documentclass{article}
\usepackage{longtable}
\usepackage[margin=1in]{geometry}
\usepackage{booktabs}
\usepackage{caption}
\usepackage{titlesec}

% Título formatado
\titleformat{\section}
  {\normalfont\Large\bfseries}{\thesection}{1em}{}
\titleformat{\subsection}
  {\normalfont\large\bfseries}{\thesubsection}{1em}{}
\titleformat{\paragraph}[runin]
  {\normalfont\normalsize\bfseries}{\theparagraph}{1em}{}

\begin{document}

\title{Análise de Entropia}
\author{}
\date{}
\maketitle

\section{Resultados da Análise}

Entropia da classe: 0.5012

\vspace{1em}

\begin{longtable}{|l|c|c|c|}
\caption{Entropia Padrão, Entropia Condicional e Ganho de Informação por Atributo} \\
\hline
\textbf{Atributo} & \textbf{Entropia Padrão} & \textbf{Entropia Condicional} & \textbf{Ganho de Informação} \\
\hline
\endfirsthead

\hline
\textbf{Atributo} & \textbf{Entropia Padrão} & \textbf{Entropia Condicional} & \textbf{Ganho de Informação} \\
\hline
\endhead

\hline
\multicolumn{4}{r}{\textit{Continua na próxima página}} \\
\endfoot

\hline
\endlastfoot

A001 & 0.7053 & 0.5008 & 0.0004 \\
A002010 & 0.0312 & 0.5012 & 0.0000 \\
A005012 & 0.5270 & 0.5010 & 0.0001 \\
A01501 & 0.6335 & 0.5012 & 0.0000 \\
C004 & 0.8905 & 0.5012 & 0.0000 \\
C006 & 0.9980 & 0.4791 & 0.0220 \\
C011 & 1.4689 & 0.4972 & 0.0040 \\
C014 & 0.9290 & 0.4988 & 0.0024 \\
D001 & 0.1504 & 0.5012 & 0.0000 \\
D00201 & 0.4019 & 0.5011 & 0.0001 \\
D00901 & 1.6610 & 0.5004 & 0.0008 \\
E011 & 0.9977 & 0.4946 & 0.0066 \\
E01401 & 2.3286 & 0.4889 & 0.0123 \\
E01403 & 0.9843 & 0.4979 & 0.0033 \\
E033 & 2.0493 & 0.4966 & 0.0046 \\
J05301 & 0.3124 & 0.4959 & 0.0052 \\
M005010 & 0.5318 & 0.5012 & 0.0000 \\
M009 & 0.9919 & 0.5012 & 0.0000 \\
M011021 & 0.7229 & 0.5008 & 0.0004 \\
M011061 & 0.2956 & 0.5012 & 0.0000 \\
M01401 & 1.1329 & 0.4989 & 0.0023 \\
M01501 & 1.2046 & 0.5001 & 0.0011 \\
M01601 & 1.4595 & 0.5002 & 0.0010 \\
M01701 & 0.9076 & 0.5009 & 0.0003 \\
M01901 & 1.5643 & 0.5003 & 0.0009 \\
N00101 & 0.8337 & 0.4664 & 0.0348 \\
N010 & 1.4284 & 0.4484 & 0.0528 \\
N011 & 1.5009 & 0.4619 & 0.0393 \\
N012 & 1.2253 & 0.4516 & 0.0496 \\
N013 & 1.0319 & 0.4569 & 0.0443 \\
N014 & 1.0156 & 0.4670 & 0.0342 \\
N015 & 0.9306 & 0.4636 & 0.0376 \\
N016 & 1.1772 & 0.4279 & 0.0733 \\
N017 & 0.8811 & 0.4523 & 0.0489 \\
N018 & 0.3251 & 0.4772 & 0.0239 \\
P03001 & 0.3258 & 0.5008 & 0.0004 \\
P03201 & 0.7731 & 0.4992 & 0.0020 \\
P03301 & 1.0305 & 0.4996 & 0.0016 \\
P03302 & 0.0822 & 0.5009 & 0.0003 \\
P03303 & 1.2246 & 0.5002 & 0.0010 \\
Q03001 & 0.2780 & 0.4989 & 0.0023 \\
Q06306 & 0.2079 & 0.4945 & 0.0067 \\
Q079 & 0.2469 & 0.4916 & 0.0096 \\
Media\_Por\_Dormitorio & 0.1945 & 0.5009 & 0.0003 \\
Salario & 2.1194 & 0.4916 & 0.0096 \\
Carga\_Horaria\_Semanal & 1.8782 & 0.4913 & 0.0099 \\
Tempo\_Deslocamento\_Min & 2.1660 & 0.4942 & 0.0070 \\
Categoria\_IMC & 1.6325 & 0.4992 & 0.0020 \\
Frequencia\_Alcool & 1.6031 & 0.4989 & 0.0023 \\
Tempo\_Minimo\_Exercicio & 0.8865 & 0.5010 & 0.0002 \\
Pratica\_Exercicio & 0.9923 & 0.5010 & 0.0002 \\
TempoDeTela & 1.0754 & 0.5010 & 0.0002 \\
Frequencia\_Fumo & 1.4303 & 0.4980 & 0.0032 \\
Violencia\_Psicologica & 0.9775 & 0.4868 & 0.0144 \\
Violencia\_Fisica & 0.3149 & 0.4975 & 0.0037 \\
Violencia\_Sexual & 0.3743 & 0.4875 & 0.0137 \\
Idade & 1.9403 & 0.4964 & 0.0048 \\
Deslocamento\_Trabalho\_Dias & 1.5432 & 0.4928 & 0.0084 \\
\end{longtable}

\subsection*{Filtragem por distribuição de valores}

Investigando os atributos com entropia padrão muito baixa, descobrimos alguns que continham distribuição desbalanceada. Criamos um corte de no mínimo  5\% para a frequência relativa do valor com menos instâncias. Alguns atributos foram preservados por terem seu valor de ganho passar por um corte de no mínimo 0.01 bits.

\paragraph{Atributos removidos:}
\texttt{A001, A002010, C011, D001, E011, E033, M01701, P03301, P03302, P03303, Q03001, Q06306, Q079, Media\_Por\_Dormitorio, Categoria\_IMC, Frequencia\_Alcool, Violencia\_Fisica}

\paragraph{Atributos preservados:}
\texttt{E01401, N00101, N012, N013, N014, N015, N016, N017, N018, Violencia\_Psicologica}

\subsection*{Filtragem por entropia condicional}

Após a filtragem anterior, analisamos a entropia condicional de cada atributo. Definimos um corte para a entropia condicional de no máximo 0.5 bits (0.012 abaixo da entropia da classe). Alguns atributos foram preservados por sua relevância contextual.

\paragraph{Atributos removidos:}
\texttt{A005012, A01501, D00201, M005010, M009, M011021, M011061, M01501, M01601, M01901, P03001, Pratica\_Exercicio}

\paragraph{Atributos preservados:}
\texttt{D00901, C004, Tempo\_Minimo\_Exercicio, TempoDeTela, Frequencia\_Fumo}

\subsection*{Filtragem empírica}

Por fim, ao analisar os 29 atributos restantes da base, observamos uma concentração de atributos da seção N, que avaliam a frequência de sintomas associados ao estado emocional nas duas últimas semanas. Para evitar \textit{overfitting} e reduzir a redundância informacional, optamos por remover os atributos \texttt{N00101}, \texttt{N011}, \texttt{N013}, \texttt{N014} e \texttt{N015}, que tratam respectivamente de como o indivíduo se sente sobre seu bem estar físico e mental, problemas por não se sentir descansado e disposto durante o dia, problemas pra se concentrar em atividades habituais, falta de apetite ou alimentação excessiva, e lentidão no movimento e fala ou agitação e inquietamento. A escolha foi feita com base no menor ganho de informação desses itens em comparação com os demais da mesma seção.

Mantivemos os atributos \texttt{N010}, \texttt{N012}, \texttt{N016} e \texttt{N017}, os quais apresentaram maior ganho e abrangem sentimentos de problemas no sono, desinteresse e falta de prazer, sentimentos de tristeza, e sentimento de fracasso ou de ser uma decepção. Apesar de \texttt{N018} (pensamentos de que seria melhor estar morto ou se machucar) apresentar ganho inferior aos atributos removidos, foi mantido devido à sua relevância e potencial importância preditiva no contexto da saúde mental.

\subsection*{Conclusão}

Após a seleção de atributos, restaram 24 variáveis na base. Os atributos preservados artificialmente apresentaram ganhos baixos (0 a 0.0008), enquanto os selecionados via filtragem apresentaram ganhos de 0.002 a 0.073, indicando maior relevância para previsão da classe.

A filtragem por distribuição identificou atributos com entropia padrão muito baixa, geralmente associados a variáveis com desbalanceamento de respostas. Entre os atributos removidos, destacam-se alguns grupos: condições de trabalho e ocupação (\texttt{E011}, \texttt{M01701}, \texttt{P03301}, \texttt{P03302}, \texttt{P03303}); variáveis relacionadas à saúde física, comportamental e emocional (\texttt{Q03001}, \texttt{Categoria\_IMC}, \texttt{Frequencia\_Alcool}, \texttt{Violencia\_Fisica}); e aspectos cognitivos e de saúde mental (\texttt{E033}, \texttt{Q06306}, \texttt{Q079}). Em menor número, foram removidas variáveis ligadas à infraestrutura domiciliar (\texttt{A001}, \texttt{A002010}), escolaridade e alfabetização básica (\texttt{D001}), condição civil ou estrutura familiar (\texttt{C011}) e indicadores derivados de densidade habitacional (\texttt{Media\_Por\_Dormitorio}).

Por outro lado, a filtragem por entropia condicional eliminou atributos que, apesar de apresentarem distribuição aceitável, não contribuíam para a redução da incerteza sobre a classe. Os atributos removidos concentram-se majoritariamente em características do trabalho e ocupação (\texttt{M005010}, \texttt{M009}, \texttt{M011021}, \texttt{M011061}, \texttt{M01501}, \texttt{M01601}, \texttt{M01901}, \texttt{P03001}), além de variáveis de infraestrutura domiciliar (\texttt{A005012}, \texttt{A01501}), escolaridade (\texttt{D00201}) e hábitos de vida (\texttt{Pratica\_Exercicio}).

\end{document}
